\section{成化至正德的补充编年节点（一）}
这一组条目把此前尚未设导航入口的年度材料接回本章脉络。它们不是把同年材料机械等同为因果，而是在可核查的时间位置上，分别保留行动、行政记录与社会经验之间的距离。

\subsection{1473年前后的材料线索}

\eventanchor{ML-1473-001}{明·1473（成化九年）}
\paragraph{1473（成化九年）：明军与蒙古盟友联合进攻哈密，但在蒙古人放弃他们后撤退} 这一节点应放回本章已经展开的权力、财政、地方社会与边地关系中理解。账本把它出现的条件概括为：前期边防、地方武装或跨区域资源调动的压力推动了该行动。 随后可追踪的变化是：它改变了后续调兵、补给或地方秩序的条件；战果和伤亡须分开核对。 可由《宪宗实录》成化九年条；《明史·兵志》及相关列传；边镇、地方志或对方材料 互相核对；不过，译写候选以编年材料定位；实录记载反映朝廷选择，崇祯期须另以奏疏、地方及敌对政权材料交叉。 因而这里标示的是材料能够支持的行动与制度位置，而不是对动机、地方执行或普通人经验的无保留复原。

\eventanchor{ML-1473-002}{明·1473（成化九年）}
\paragraph{1473（成化九年）：科举、学校和法律条文在本年制度材料中可追踪，不能据规范文本推定州县实践一致。} 这一节点应放回本章已经展开的权力、财政、地方社会与边地关系中理解。账本把它出现的条件概括为：制度、教育或知识传播的需要推动了相关行动或文本形成。 随后可追踪的变化是：它提供制度和传播史的锚点，不能据此推断社会整体接受。 可由《宪宗实录》成化九年条；《明史·选举志》或《艺文志》；文集、碑刻或书目材料 互相核对；不过，桥接条目只标示可回查的年度问题与材料链；实录有中央选择性，崇祯期尤其需要奏疏、地方、朝鲜及满文材料交叉。 因而这里标示的是材料能够支持的行动与制度位置，而不是对动机、地方执行或普通人经验的无保留复原。

\eventanchor{ML-1473-003}{明·1473（成化九年）}
\paragraph{1473（成化九年）：土木危机后的边镇整顿、京师防务或复辟前后军政调度在本年材料中构成连续问题。} 这一节点应放回本章已经展开的权力、财政、地方社会与边地关系中理解。账本把它出现的条件概括为：前期边防、地方武装或跨区域资源调动的压力推动了该行动。 随后可追踪的变化是：它改变了后续调兵、补给或地方秩序的条件；战果和伤亡须分开核对。 可由《宪宗实录》成化九年条；《明史·兵志》及相关列传；边镇、地方志或对方材料 互相核对；不过，桥接条目只标示可回查的年度问题与材料链；实录有中央选择性，崇祯期尤其需要奏疏、地方、朝鲜及满文材料交叉。 因而这里标示的是材料能够支持的行动与制度位置，而不是对动机、地方执行或普通人经验的无保留复原。

\eventanchor{MM-1473-HAMI-SUCCESSION}{明·1473（成化九年）}
\paragraph{1473（成化九年）：哈密王位继承与朝贡使团问题引发明廷交涉，吐鲁番势力介入东天山政治} 这一节点应放回本章已经展开的权力、财政、地方社会与边地关系中理解。账本把它出现的条件概括为：哈密既是明朝西域朝贡体系的节点，也受本地王族和吐鲁番力量竞争牵动。 随后可追踪的变化是：明廷不得不在册封、赏赐、关市与边防之间权衡；哈密问题此后反复成为西北边政负担。 可由《宪宗实录》成化九年哈密条；《明史·西域传》；《明实录》所载哈密使臣文书 互相核对；不过，哈密政局和明廷交涉可见于实录；当地王位传承、部众立场和吐鲁番影响范围在中原材料中较难确证。 因而这里标示的是材料能够支持的行动与制度位置，而不是对动机、地方执行或普通人经验的无保留复原。

\eventanchor{ML-1474-001}{明·1474（成化十年）}
\paragraph{1474（成化十年）：于子军指挥重建和扩建长城，南封鄂尔多斯} 这一节点应放回本章已经展开的权力、财政、地方社会与边地关系中理解。账本把它出现的条件概括为：继承、官僚协调或皇权运作的压力使问题进入朝廷程序。 随后可追踪的变化是：它影响后续人事与政策议程；动机和派系标签不能由单一叙事裁定。 可由《宪宗实录》成化十年条；《明史·本纪》及相关列传；诏令、奏疏或会典版本 互相核对；不过，译写候选以编年材料定位；实录记载反映朝廷选择，崇祯期须另以奏疏、地方及敌对政权材料交叉。 因而这里标示的是材料能够支持的行动与制度位置，而不是对动机、地方执行或普通人经验的无保留复原。

\eventanchor{ML-1474-002}{明·1474（成化十年）}
\paragraph{1474（成化十年）：科举、学校和法律条文在本年制度材料中可追踪，不能据规范文本推定州县实践一致。} 这一节点应放回本章已经展开的权力、财政、地方社会与边地关系中理解。账本把它出现的条件概括为：制度、教育或知识传播的需要推动了相关行动或文本形成。 随后可追踪的变化是：它提供制度和传播史的锚点，不能据此推断社会整体接受。 可由《宪宗实录》成化十年条；《明史·选举志》或《艺文志》；文集、碑刻或书目材料 互相核对；不过，桥接条目只标示可回查的年度问题与材料链；实录有中央选择性，崇祯期尤其需要奏疏、地方、朝鲜及满文材料交叉。 因而这里标示的是材料能够支持的行动与制度位置，而不是对动机、地方执行或普通人经验的无保留复原。

\eventanchor{ML-1474-003}{明·1474（成化十年）}
\paragraph{1474（成化十年）：地方灾异、赈济与水利条目为本年社会史提供入口，地方志的修纂层次需要说明。} 这一节点应放回本章已经展开的权力、财政、地方社会与边地关系中理解。账本把它出现的条件概括为：自然风险与地方报告促成朝廷或地方的应对记录。 随后可追踪的变化是：赈济、迁徙和损失的实际范围仍须以受灾地区材料分别核验。 可由《宪宗实录》成化十年条；《明史·五行志》；受灾府县地方志与奏疏 互相核对；不过，桥接条目只标示可回查的年度问题与材料链；实录有中央选择性，崇祯期尤其需要奏疏、地方、朝鲜及满文材料交叉。 因而这里标示的是材料能够支持的行动与制度位置，而不是对动机、地方执行或普通人经验的无保留复原。

\eventanchor{ML-1474-004}{明·1474（成化十年）}
\paragraph{1474（成化十年）：景泰、天顺时期的继承、人事与诏令链条可在本年材料中定位，政变责任不可由单一史书裁断。} 这一节点应放回本章已经展开的权力、财政、地方社会与边地关系中理解。账本把它出现的条件概括为：继承、官僚协调或皇权运作的压力使问题进入朝廷程序。 随后可追踪的变化是：它影响后续人事与政策议程；动机和派系标签不能由单一叙事裁定。 可由《宪宗实录》成化十年条；《明史·本纪》及相关列传；诏令、奏疏或会典版本 互相核对；不过，桥接条目只标示可回查的年度问题与材料链；实录有中央选择性，崇祯期尤其需要奏疏、地方、朝鲜及满文材料交叉。 因而这里标示的是材料能够支持的行动与制度位置，而不是对动机、地方执行或普通人经验的无保留复原。

\eventanchor{ML-1475-001}{明·1475（成化十一年）}
\paragraph{1475（成化十一年）：湖南苗族叛乱遭镇压} 这一节点应放回本章已经展开的权力、财政、地方社会与边地关系中理解。账本把它出现的条件概括为：继承、官僚协调或皇权运作的压力使问题进入朝廷程序。 随后可追踪的变化是：它影响后续人事与政策议程；动机和派系标签不能由单一叙事裁定。 可由《宪宗实录》成化十一年条；《明史·本纪》及相关列传；诏令、奏疏或会典版本 互相核对；不过，译写候选以编年材料定位；实录记载反映朝廷选择，崇祯期须另以奏疏、地方及敌对政权材料交叉。 因而这里标示的是材料能够支持的行动与制度位置，而不是对动机、地方执行或普通人经验的无保留复原。

\eventanchor{ML-1475-002}{明·1475（成化十一年）}
\paragraph{1475（成化十一年）：学校、考试、刻书或礼制材料在本年留下文化制度线索，精英文本不能代表全体社会。} 这一节点应放回本章已经展开的权力、财政、地方社会与边地关系中理解。账本把它出现的条件概括为：制度、教育或知识传播的需要推动了相关行动或文本形成。 随后可追踪的变化是：它提供制度和传播史的锚点，不能据此推断社会整体接受。 可由《宪宗实录》成化十一年条；《明史·选举志》或《艺文志》；文集、碑刻或书目材料 互相核对；不过，桥接条目只标示可回查的年度问题与材料链；实录有中央选择性，崇祯期尤其需要奏疏、地方、朝鲜及满文材料交叉。 因而这里标示的是材料能够支持的行动与制度位置，而不是对动机、地方执行或普通人经验的无保留复原。

\eventanchor{ML-1475-003}{明·1475（成化十一年）}
\paragraph{1475（成化十一年）：本年灾荒、河工或赈恤的报奏材料可与地方志互校，报灾范围和实际损失应分别记录。（材料核查重点：诏令或奏议的形成与发受链）} 这一节点应放回本章已经展开的权力、财政、地方社会与边地关系中理解。账本把它出现的条件概括为：自然风险与地方报告促成朝廷或地方的应对记录。 随后可追踪的变化是：赈济、迁徙和损失的实际范围仍须以受灾地区材料分别核验。 可由《宪宗实录》成化十一年条；《明史·五行志》；受灾府县地方志与奏疏 互相核对；不过，桥接条目只标示可回查的年度问题与材料链；实录有中央选择性，崇祯期尤其需要奏疏、地方、朝鲜及满文材料交叉。；本条特别核对诏令或奏议的形成与发受链。 因而这里标示的是材料能够支持的行动与制度位置，而不是对动机、地方执行或普通人经验的无保留复原。

\eventanchor{ML-1475-004}{明·1475（成化十一年）}
\paragraph{1475（成化十一年）：本年灾荒、河工或赈恤的报奏材料可与地方志互校，报灾范围和实际损失应分别记录。（材料核查重点：报称信息的来源、抄传与行政分类）} 这一节点应放回本章已经展开的权力、财政、地方社会与边地关系中理解。账本把它出现的条件概括为：自然风险与地方报告促成朝廷或地方的应对记录。 随后可追踪的变化是：赈济、迁徙和损失的实际范围仍须以受灾地区材料分别核验。 可由《宪宗实录》成化十一年条；《明史·五行志》；受灾府县地方志与奏疏 互相核对；不过，桥接条目只标示可回查的年度问题与材料链；实录有中央选择性，崇祯期尤其需要奏疏、地方、朝鲜及满文材料交叉。；本条特别核对报称信息的来源、抄传与行政分类。 因而这里标示的是材料能够支持的行动与制度位置，而不是对动机、地方执行或普通人经验的无保留复原。

\eventanchor{ML-1475-005}{明·1475（成化十一年）}
\paragraph{1475（成化十一年）：本年灾荒、河工或赈恤的报奏材料可与地方志互校，报灾范围和实际损失应分别记录。（材料核查重点：同年地方材料所见的执行范围）} 这一节点应放回本章已经展开的权力、财政、地方社会与边地关系中理解。账本把它出现的条件概括为：自然风险与地方报告促成朝廷或地方的应对记录。 随后可追踪的变化是：赈济、迁徙和损失的实际范围仍须以受灾地区材料分别核验。 可由《宪宗实录》成化十一年条；《明史·五行志》；受灾府县地方志与奏疏 互相核对；不过，桥接条目只标示可回查的年度问题与材料链；实录有中央选择性，崇祯期尤其需要奏疏、地方、朝鲜及满文材料交叉。；本条特别核对同年地方材料所见的执行范围。 因而这里标示的是材料能够支持的行动与制度位置，而不是对动机、地方执行或普通人经验的无保留复原。

\eventanchor{ML-1476-001}{明·1476（成化十二年）六月}
\paragraph{1476（成化十二年）六月：襄阳周边的流浪人口叛乱，直到政府允许他们减税索取土地} 这一节点应放回本章已经展开的权力、财政、地方社会与边地关系中理解。账本把它出现的条件概括为：财政征解、交通工程或货币支付的压力使相关措施进入政务。 随后可追踪的变化是：其法定安排不等于地方执行，后续须以地域材料追踪负担与成效。 可由《宪宗实录》成化十二年条；《明会典》相关条目；地方志、黄册/契约/河工材料 互相核对；不过，译写候选以编年材料定位；实录记载反映朝廷选择，崇祯期须另以奏疏、地方及敌对政权材料交叉。 因而这里标示的是材料能够支持的行动与制度位置，而不是对动机、地方执行或普通人经验的无保留复原。

\eventanchor{ML-1493-002}{明·1493（弘治六年）}
\paragraph{1493（弘治六年）：学校、考试、刻书或礼制材料在本年留下文化制度线索，精英文本不能代表全体社会。（材料核查重点：诏令或奏议的形成与发受链）} 这一节点应放回本章已经展开的权力、财政、地方社会与边地关系中理解。账本把它出现的条件概括为：制度、教育或知识传播的需要推动了相关行动或文本形成。 随后可追踪的变化是：它提供制度和传播史的锚点，不能据此推断社会整体接受。 可由《孝宗实录》弘治六年条；《明史·选举志》或《艺文志》；文集、碑刻或书目材料 互相核对；不过，桥接条目只标示可回查的年度问题与材料链；实录有中央选择性，崇祯期尤其需要奏疏、地方、朝鲜及满文材料交叉。；本条特别核对诏令或奏议的形成与发受链。 因而这里标示的是材料能够支持的行动与制度位置，而不是对动机、地方执行或普通人经验的无保留复原。

\eventanchor{ML-1493-003}{明·1493（弘治六年）}
\paragraph{1493（弘治六年）：学校、考试、刻书或礼制材料在本年留下文化制度线索，精英文本不能代表全体社会。（材料核查重点：报称信息的来源、抄传与行政分类）} 这一节点应放回本章已经展开的权力、财政、地方社会与边地关系中理解。账本把它出现的条件概括为：制度、教育或知识传播的需要推动了相关行动或文本形成。 随后可追踪的变化是：它提供制度和传播史的锚点，不能据此推断社会整体接受。 可由《孝宗实录》弘治六年条；《明史·选举志》或《艺文志》；文集、碑刻或书目材料 互相核对；不过，桥接条目只标示可回查的年度问题与材料链；实录有中央选择性，崇祯期尤其需要奏疏、地方、朝鲜及满文材料交叉。；本条特别核对报称信息的来源、抄传与行政分类。 因而这里标示的是材料能够支持的行动与制度位置，而不是对动机、地方执行或普通人经验的无保留复原。

\eventanchor{MM-1493-TURPAN-HAMI-TAKEOVER}{明·1493（弘治六年）}
\paragraph{1493（弘治六年）：吐鲁番军进入哈密并改置亲近的统治者，明朝失去对该绿洲的直接影响} 这一节点应放回本章已经展开的权力、财政、地方社会与边地关系中理解。账本把它出现的条件概括为：前次册封未能建立可自保的哈密政权，吐鲁番以邻近兵力和王族关系扩大在绿洲的影响。 随后可追踪的变化是：明廷以限制入贡和重新扶持继承人为回应；嘉峪关外的贸易、使团和情报网络更不稳定。 可由《孝宗实录》弘治六年哈密条；《明史·西域传》；嘉峪关边官奏疏 互相核对；不过，吐鲁番夺取哈密及明廷交涉可由实录和《明史》互见；行动路线、死伤和哈密居民反应多经使臣与边官转述。 因而这里标示的是材料能够支持的行动与制度位置，而不是对动机、地方执行或普通人经验的无保留复原。

\eventanchor{ML-1494-001}{明·1494（弘治七年）}
\paragraph{1494（弘治七年）：1494年黄河洪水：黄河洪水泛滥，但刘大侠成功引导黄河流向山东以南，稳定了黄河的航向，直到19世纪} 这一节点应放回本章已经展开的权力、财政、地方社会与边地关系中理解。账本把它出现的条件概括为：自然风险与地方报告促成朝廷或地方的应对记录。 随后可追踪的变化是：赈济、迁徙和损失的实际范围仍须以受灾地区材料分别核验。 可由《孝宗实录》弘治七年条；《明史·五行志》；受灾府县地方志与奏疏 互相核对；不过，译写候选以编年材料定位；实录记载反映朝廷选择，崇祯期须另以奏疏、地方及敌对政权材料交叉。 因而这里标示的是材料能够支持的行动与制度位置，而不是对动机、地方执行或普通人经验的无保留复原。

\eventanchor{ML-1494-002}{明·1494（弘治七年）}
\paragraph{1494（弘治七年）：国家军队改革转向为地方部队招募志愿者} 这一节点应放回本章已经展开的权力、财政、地方社会与边地关系中理解。账本把它出现的条件概括为：前期边防、地方武装或跨区域资源调动的压力推动了该行动。 随后可追踪的变化是：它改变了后续调兵、补给或地方秩序的条件；战果和伤亡须分开核对。 可由《孝宗实录》弘治七年条；《明史·兵志》及相关列传；边镇、地方志或对方材料 互相核对；不过，译写候选以编年材料定位；实录记载反映朝廷选择，崇祯期须另以奏疏、地方及敌对政权材料交叉。 因而这里标示的是材料能够支持的行动与制度位置，而不是对动机、地方执行或普通人经验的无保留复原。

\eventanchor{ML-1494-003}{明·1494（弘治七年）}
\paragraph{1494（弘治七年）：本年灾荒、河工或赈恤的报奏材料可与地方志互校，报灾范围和实际损失应分别记录。（材料核查重点：诏令或奏议的形成与发受链）} 这一节点应放回本章已经展开的权力、财政、地方社会与边地关系中理解。账本把它出现的条件概括为：自然风险与地方报告促成朝廷或地方的应对记录。 随后可追踪的变化是：赈济、迁徙和损失的实际范围仍须以受灾地区材料分别核验。 可由《孝宗实录》弘治七年条；《明史·五行志》；受灾府县地方志与奏疏 互相核对；不过，桥接条目只标示可回查的年度问题与材料链；实录有中央选择性，崇祯期尤其需要奏疏、地方、朝鲜及满文材料交叉。；本条特别核对诏令或奏议的形成与发受链。 因而这里标示的是材料能够支持的行动与制度位置，而不是对动机、地方执行或普通人经验的无保留复原。

\subsection{1494年前后的材料线索}

\eventanchor{ML-1494-004}{明·1494（弘治七年）}
\paragraph{1494（弘治七年）：本年灾荒、河工或赈恤的报奏材料可与地方志互校，报灾范围和实际损失应分别记录。（材料核查重点：报称信息的来源、抄传与行政分类）} 这一节点应放回本章已经展开的权力、财政、地方社会与边地关系中理解。账本把它出现的条件概括为：自然风险与地方报告促成朝廷或地方的应对记录。 随后可追踪的变化是：赈济、迁徙和损失的实际范围仍须以受灾地区材料分别核验。 可由《孝宗实录》弘治七年条；《明史·五行志》；受灾府县地方志与奏疏 互相核对；不过，桥接条目只标示可回查的年度问题与材料链；实录有中央选择性，崇祯期尤其需要奏疏、地方、朝鲜及满文材料交叉。；本条特别核对报称信息的来源、抄传与行政分类。 因而这里标示的是材料能够支持的行动与制度位置，而不是对动机、地方执行或普通人经验的无保留复原。

\eventanchor{ML-1495-001}{明·1495（弘治八年）}
\paragraph{1495（弘治八年）：明军短暂占领哈密，随后来自吐鲁番的增援迫使他们撤退} 这一节点应放回本章已经展开的权力、财政、地方社会与边地关系中理解。账本把它出现的条件概括为：前期边防、地方武装或跨区域资源调动的压力推动了该行动。 随后可追踪的变化是：它改变了后续调兵、补给或地方秩序的条件；战果和伤亡须分开核对。 可由《孝宗实录》弘治八年条；《明史·兵志》及相关列传；边镇、地方志或对方材料 互相核对；不过，译写候选以编年材料定位；实录记载反映朝廷选择，崇祯期须另以奏疏、地方及敌对政权材料交叉。 因而这里标示的是材料能够支持的行动与制度位置，而不是对动机、地方执行或普通人经验的无保留复原。

\eventanchor{ML-1495-002}{明·1495（弘治八年）}
\paragraph{1495（弘治八年）：成化、弘治间中枢任免、厂卫与言路的本年材料标记皇权和官僚关系的调整。（材料核查重点：诏令或奏议的形成与发受链）} 这一节点应放回本章已经展开的权力、财政、地方社会与边地关系中理解。账本把它出现的条件概括为：继承、官僚协调或皇权运作的压力使问题进入朝廷程序。 随后可追踪的变化是：它影响后续人事与政策议程；动机和派系标签不能由单一叙事裁定。 可由《孝宗实录》弘治八年条；《明史·本纪》及相关列传；诏令、奏疏或会典版本 互相核对；不过，桥接条目只标示可回查的年度问题与材料链；实录有中央选择性，崇祯期尤其需要奏疏、地方、朝鲜及满文材料交叉。；本条特别核对诏令或奏议的形成与发受链。 因而这里标示的是材料能够支持的行动与制度位置，而不是对动机、地方执行或普通人经验的无保留复原。

\eventanchor{ML-1495-003}{明·1495（弘治八年）}
\paragraph{1495（弘治八年）：成化、弘治间中枢任免、厂卫与言路的本年材料标记皇权和官僚关系的调整。（材料核查重点：报称信息的来源、抄传与行政分类）} 这一节点应放回本章已经展开的权力、财政、地方社会与边地关系中理解。账本把它出现的条件概括为：继承、官僚协调或皇权运作的压力使问题进入朝廷程序。 随后可追踪的变化是：它影响后续人事与政策议程；动机和派系标签不能由单一叙事裁定。 可由《孝宗实录》弘治八年条；《明史·本纪》及相关列传；诏令、奏疏或会典版本 互相核对；不过，桥接条目只标示可回查的年度问题与材料链；实录有中央选择性，崇祯期尤其需要奏疏、地方、朝鲜及满文材料交叉。；本条特别核对报称信息的来源、抄传与行政分类。 因而这里标示的是材料能够支持的行动与制度位置，而不是对动机、地方执行或普通人经验的无保留复原。

\eventanchor{MM-1495-HAMI-RESTORATION}{明·1495（弘治八年）}
\paragraph{1495（弘治八年）：获明廷承认的哈密继承人恢复入贡，哈密一度脱离吐鲁番的直接控制} 这一节点应放回本章已经展开的权力、财政、地方社会与边地关系中理解。账本把它出现的条件概括为：吐鲁番占领引发哈密王族与商路相关群体的不满，明廷也希望保留嘉峪关外的缓冲节点。 随后可追踪的变化是：朝贡秩序暂获修复，但哈密仍须在吐鲁番军力和明朝有限支援之间周旋；这并非稳定复国。 可由《孝宗实录》弘治八年哈密条；《明史·西域传》；哈密使臣入贡文书 互相核对；不过，恢复朝贡和册封关系可证；明廷援助的军事实效、恢复范围和持续时间须区别使团文书与后出边疆叙事。 因而这里标示的是材料能够支持的行动与制度位置，而不是对动机、地方执行或普通人经验的无保留复原。

\eventanchor{ML-1496-001}{明·1496（弘治九年）}
\paragraph{1496（弘治九年）：日本使节到明朝中国：日本特使从北京回程时杀害数人} 这一节点应放回本章已经展开的权力、财政、地方社会与边地关系中理解。账本把它出现的条件概括为：朝贡名义、边境安全与港口贸易的利益使交涉持续发生。 随后可追踪的变化是：它为后续册封、互市或海防安排留下文书与跨政权比较线索。 可由《孝宗实录》弘治九年条；《明史·外国传》；《朝鲜王朝实录》、琉球《历代宝案》或海外档案 互相核对；不过，译写候选以编年材料定位；实录记载反映朝廷选择，崇祯期须另以奏疏、地方及敌对政权材料交叉。 因而这里标示的是材料能够支持的行动与制度位置，而不是对动机、地方执行或普通人经验的无保留复原。

\eventanchor{ML-1496-002}{明·1496（弘治九年）}
\paragraph{1496（弘治九年）：田土、赋役、盐课或漕运的本年政策记录提供财政执行的候选节点，需同地区册籍比照。} 这一节点应放回本章已经展开的权力、财政、地方社会与边地关系中理解。账本把它出现的条件概括为：财政征解、交通工程或货币支付的压力使相关措施进入政务。 随后可追踪的变化是：其法定安排不等于地方执行，后续须以地域材料追踪负担与成效。 可由《孝宗实录》弘治九年条；《明会典》相关条目；地方志、黄册/契约/河工材料 互相核对；不过，桥接条目只标示可回查的年度问题与材料链；实录有中央选择性，崇祯期尤其需要奏疏、地方、朝鲜及满文材料交叉。 因而这里标示的是材料能够支持的行动与制度位置，而不是对动机、地方执行或普通人经验的无保留复原。

\eventanchor{ML-1496-003}{明·1496（弘治九年）}
\paragraph{1496（弘治九年）：荆襄、广西、辽东或西北的兵事和安置文书构成本年地方秩序重组的线索。（材料核查重点：诏令或奏议的形成与发受链）} 这一节点应放回本章已经展开的权力、财政、地方社会与边地关系中理解。账本把它出现的条件概括为：前期边防、地方武装或跨区域资源调动的压力推动了该行动。 随后可追踪的变化是：它改变了后续调兵、补给或地方秩序的条件；战果和伤亡须分开核对。 可由《孝宗实录》弘治九年条；《明史·兵志》及相关列传；边镇、地方志或对方材料 互相核对；不过，桥接条目只标示可回查的年度问题与材料链；实录有中央选择性，崇祯期尤其需要奏疏、地方、朝鲜及满文材料交叉。；本条特别核对诏令或奏议的形成与发受链。 因而这里标示的是材料能够支持的行动与制度位置，而不是对动机、地方执行或普通人经验的无保留复原。

\eventanchor{ML-1496-004}{明·1496（弘治九年）}
\paragraph{1496（弘治九年）：荆襄、广西、辽东或西北的兵事和安置文书构成本年地方秩序重组的线索。（材料核查重点：报称信息的来源、抄传与行政分类）} 这一节点应放回本章已经展开的权力、财政、地方社会与边地关系中理解。账本把它出现的条件概括为：前期边防、地方武装或跨区域资源调动的压力推动了该行动。 随后可追踪的变化是：它改变了后续调兵、补给或地方秩序的条件；战果和伤亡须分开核对。 可由《孝宗实录》弘治九年条；《明史·兵志》及相关列传；边镇、地方志或对方材料 互相核对；不过，桥接条目只标示可回查的年度问题与材料链；实录有中央选择性，崇祯期尤其需要奏疏、地方、朝鲜及满文材料交叉。；本条特别核对报称信息的来源、抄传与行政分类。 因而这里标示的是材料能够支持的行动与制度位置，而不是对动机、地方执行或普通人经验的无保留复原。

\eventanchor{ML-1497-001}{明·1497（弘治十年）}
\paragraph{1497（弘治十年）：田土、赋役、盐课或漕运的本年政策记录提供财政执行的候选节点，需同地区册籍比照。（材料核查重点：诏令或奏议的形成与发受链）} 这一节点应放回本章已经展开的权力、财政、地方社会与边地关系中理解。账本把它出现的条件概括为：财政征解、交通工程或货币支付的压力使相关措施进入政务。 随后可追踪的变化是：其法定安排不等于地方执行，后续须以地域材料追踪负担与成效。 可由《孝宗实录》弘治十年条；《明会典》相关条目；地方志、黄册/契约/河工材料 互相核对；不过，桥接条目只标示可回查的年度问题与材料链；实录有中央选择性，崇祯期尤其需要奏疏、地方、朝鲜及满文材料交叉。；本条特别核对诏令或奏议的形成与发受链。 因而这里标示的是材料能够支持的行动与制度位置，而不是对动机、地方执行或普通人经验的无保留复原。

\eventanchor{ML-1497-002}{明·1497（弘治十年）}
\paragraph{1497（弘治十年）：哈密、吐鲁番及周边朝贡关系的本年交涉显示东天山政治流动，册封不等于稳定控制。} 这一节点应放回本章已经展开的权力、财政、地方社会与边地关系中理解。账本把它出现的条件概括为：朝贡名义、边境安全与港口贸易的利益使交涉持续发生。 随后可追踪的变化是：它为后续册封、互市或海防安排留下文书与跨政权比较线索。 可由《孝宗实录》弘治十年条；《明史·外国传》；《朝鲜王朝实录》、琉球《历代宝案》或海外档案 互相核对；不过，桥接条目只标示可回查的年度问题与材料链；实录有中央选择性，崇祯期尤其需要奏疏、地方、朝鲜及满文材料交叉。 因而这里标示的是材料能够支持的行动与制度位置，而不是对动机、地方执行或普通人经验的无保留复原。

\eventanchor{ML-1497-003}{明·1497（弘治十年）}
\paragraph{1497（弘治十年）：田土、赋役、盐课或漕运的本年政策记录提供财政执行的候选节点，需同地区册籍比照。（材料核查重点：报称信息的来源、抄传与行政分类）} 这一节点应放回本章已经展开的权力、财政、地方社会与边地关系中理解。账本把它出现的条件概括为：财政征解、交通工程或货币支付的压力使相关措施进入政务。 随后可追踪的变化是：其法定安排不等于地方执行，后续须以地域材料追踪负担与成效。 可由《孝宗实录》弘治十年条；《明会典》相关条目；地方志、黄册/契约/河工材料 互相核对；不过，桥接条目只标示可回查的年度问题与材料链；实录有中央选择性，崇祯期尤其需要奏疏、地方、朝鲜及满文材料交叉。；本条特别核对报称信息的来源、抄传与行政分类。 因而这里标示的是材料能够支持的行动与制度位置，而不是对动机、地方执行或普通人经验的无保留复原。

\eventanchor{ML-1497-004}{明·1497（弘治十年）}
\paragraph{1497（弘治十年）：田土、赋役、盐课或漕运的本年政策记录提供财政执行的候选节点，需同地区册籍比照。（材料核查重点：同年地方材料所见的执行范围）} 这一节点应放回本章已经展开的权力、财政、地方社会与边地关系中理解。账本把它出现的条件概括为：财政征解、交通工程或货币支付的压力使相关措施进入政务。 随后可追踪的变化是：其法定安排不等于地方执行，后续须以地域材料追踪负担与成效。 可由《孝宗实录》弘治十年条；《明会典》相关条目；地方志、黄册/契约/河工材料 互相核对；不过，桥接条目只标示可回查的年度问题与材料链；实录有中央选择性，崇祯期尤其需要奏疏、地方、朝鲜及满文材料交叉。；本条特别核对同年地方材料所见的执行范围。 因而这里标示的是材料能够支持的行动与制度位置，而不是对动机、地方执行或普通人经验的无保留复原。

\eventanchor{ML-1498-001}{明·1498（弘治十一年）}
\paragraph{1498（弘治十一年）：哈密、吐鲁番及周边朝贡关系的本年交涉显示东天山政治流动，册封不等于稳定控制。（材料核查重点：诏令或奏议的形成与发受链）} 这一节点应放回本章已经展开的权力、财政、地方社会与边地关系中理解。账本把它出现的条件概括为：朝贡名义、边境安全与港口贸易的利益使交涉持续发生。 随后可追踪的变化是：它为后续册封、互市或海防安排留下文书与跨政权比较线索。 可由《孝宗实录》弘治十一年条；《明史·外国传》；《朝鲜王朝实录》、琉球《历代宝案》或海外档案 互相核对；不过，桥接条目只标示可回查的年度问题与材料链；实录有中央选择性，崇祯期尤其需要奏疏、地方、朝鲜及满文材料交叉。；本条特别核对诏令或奏议的形成与发受链。 因而这里标示的是材料能够支持的行动与制度位置，而不是对动机、地方执行或普通人经验的无保留复原。

\eventanchor{ML-1498-002}{明·1498（弘治十一年）}
\paragraph{1498（弘治十一年）：学校、考试、刻书或礼制材料在本年留下文化制度线索，精英文本不能代表全体社会。} 这一节点应放回本章已经展开的权力、财政、地方社会与边地关系中理解。账本把它出现的条件概括为：制度、教育或知识传播的需要推动了相关行动或文本形成。 随后可追踪的变化是：它提供制度和传播史的锚点，不能据此推断社会整体接受。 可由《孝宗实录》弘治十一年条；《明史·选举志》或《艺文志》；文集、碑刻或书目材料 互相核对；不过，桥接条目只标示可回查的年度问题与材料链；实录有中央选择性，崇祯期尤其需要奏疏、地方、朝鲜及满文材料交叉。 因而这里标示的是材料能够支持的行动与制度位置，而不是对动机、地方执行或普通人经验的无保留复原。

\eventanchor{ML-1498-003}{明·1498（弘治十一年）}
\paragraph{1498（弘治十一年）：哈密、吐鲁番及周边朝贡关系的本年交涉显示东天山政治流动，册封不等于稳定控制。（材料核查重点：报称信息的来源、抄传与行政分类）} 这一节点应放回本章已经展开的权力、财政、地方社会与边地关系中理解。账本把它出现的条件概括为：朝贡名义、边境安全与港口贸易的利益使交涉持续发生。 随后可追踪的变化是：它为后续册封、互市或海防安排留下文书与跨政权比较线索。 可由《孝宗实录》弘治十一年条；《明史·外国传》；《朝鲜王朝实录》、琉球《历代宝案》或海外档案 互相核对；不过，桥接条目只标示可回查的年度问题与材料链；实录有中央选择性，崇祯期尤其需要奏疏、地方、朝鲜及满文材料交叉。；本条特别核对报称信息的来源、抄传与行政分类。 因而这里标示的是材料能够支持的行动与制度位置，而不是对动机、地方执行或普通人经验的无保留复原。

\eventanchor{ML-1498-004}{明·1498（弘治十一年）}
\paragraph{1498（弘治十一年）：哈密、吐鲁番及周边朝贡关系的本年交涉显示东天山政治流动，册封不等于稳定控制。（材料核查重点：同年地方材料所见的执行范围）} 这一节点应放回本章已经展开的权力、财政、地方社会与边地关系中理解。账本把它出现的条件概括为：朝贡名义、边境安全与港口贸易的利益使交涉持续发生。 随后可追踪的变化是：它为后续册封、互市或海防安排留下文书与跨政权比较线索。 可由《孝宗实录》弘治十一年条；《明史·外国传》；《朝鲜王朝实录》、琉球《历代宝案》或海外档案 互相核对；不过，桥接条目只标示可回查的年度问题与材料链；实录有中央选择性，崇祯期尤其需要奏疏、地方、朝鲜及满文材料交叉。；本条特别核对同年地方材料所见的执行范围。 因而这里标示的是材料能够支持的行动与制度位置，而不是对动机、地方执行或普通人经验的无保留复原。

\eventanchor{ML-1499-001}{明·1499（弘治十二年）}
\paragraph{1499（弘治十二年）：对吐鲁番的贸易禁运迫使他们将哈密交还给维吾尔人控制} 这一节点应放回本章已经展开的权力、财政、地方社会与边地关系中理解。账本把它出现的条件概括为：财政征解、交通工程或货币支付的压力使相关措施进入政务。 随后可追踪的变化是：其法定安排不等于地方执行，后续须以地域材料追踪负担与成效。 可由《孝宗实录》弘治十二年条；《明会典》相关条目；地方志、黄册/契约/河工材料 互相核对；不过，译写候选以编年材料定位；实录记载反映朝廷选择，崇祯期须另以奏疏、地方及敌对政权材料交叉。 因而这里标示的是材料能够支持的行动与制度位置，而不是对动机、地方执行或普通人经验的无保留复原。

\eventanchor{ML-1499-002}{明·1499（弘治十二年）}
\paragraph{1499（弘治十二年）：贵州彝族起义} 这一节点应放回本章已经展开的权力、财政、地方社会与边地关系中理解。账本把它出现的条件概括为：继承、官僚协调或皇权运作的压力使问题进入朝廷程序。 随后可追踪的变化是：它影响后续人事与政策议程；动机和派系标签不能由单一叙事裁定。 可由《孝宗实录》弘治十二年条；《明史·本纪》及相关列传；诏令、奏疏或会典版本 互相核对；不过，译写候选以编年材料定位；实录记载反映朝廷选择，崇祯期须另以奏疏、地方及敌对政权材料交叉。 因而这里标示的是材料能够支持的行动与制度位置，而不是对动机、地方执行或普通人经验的无保留复原。

\eventanchor{ML-1499-003}{明·1499（弘治十二年）}
\paragraph{1499（弘治十二年）：学校、考试、刻书或礼制材料在本年留下文化制度线索，精英文本不能代表全体社会。（材料核查重点：诏令或奏议的形成与发受链）} 这一节点应放回本章已经展开的权力、财政、地方社会与边地关系中理解。账本把它出现的条件概括为：制度、教育或知识传播的需要推动了相关行动或文本形成。 随后可追踪的变化是：它提供制度和传播史的锚点，不能据此推断社会整体接受。 可由《孝宗实录》弘治十二年条；《明史·选举志》或《艺文志》；文集、碑刻或书目材料 互相核对；不过，桥接条目只标示可回查的年度问题与材料链；实录有中央选择性，崇祯期尤其需要奏疏、地方、朝鲜及满文材料交叉。；本条特别核对诏令或奏议的形成与发受链。 因而这里标示的是材料能够支持的行动与制度位置，而不是对动机、地方执行或普通人经验的无保留复原。

\subsection{1499年前后的材料线索}

\eventanchor{ML-1499-004}{明·1499（弘治十二年）}
\paragraph{1499（弘治十二年）：学校、考试、刻书或礼制材料在本年留下文化制度线索，精英文本不能代表全体社会。（材料核查重点：报称信息的来源、抄传与行政分类）} 这一节点应放回本章已经展开的权力、财政、地方社会与边地关系中理解。账本把它出现的条件概括为：制度、教育或知识传播的需要推动了相关行动或文本形成。 随后可追踪的变化是：它提供制度和传播史的锚点，不能据此推断社会整体接受。 可由《孝宗实录》弘治十二年条；《明史·选举志》或《艺文志》；文集、碑刻或书目材料 互相核对；不过，桥接条目只标示可回查的年度问题与材料链；实录有中央选择性，崇祯期尤其需要奏疏、地方、朝鲜及满文材料交叉。；本条特别核对报称信息的来源、抄传与行政分类。 因而这里标示的是材料能够支持的行动与制度位置，而不是对动机、地方执行或普通人经验的无保留复原。

\eventanchor{ML-1500-001}{明·1500（弘治十三年）}
\paragraph{1500（弘治十三年）：海南黎族起义} 这一节点应放回本章已经展开的权力、财政、地方社会与边地关系中理解。账本把它出现的条件概括为：继承、官僚协调或皇权运作的压力使问题进入朝廷程序。 随后可追踪的变化是：它影响后续人事与政策议程；动机和派系标签不能由单一叙事裁定。 可由《孝宗实录》弘治十三年条；《明史·本纪》及相关列传；诏令、奏疏或会典版本 互相核对；不过，译写候选以编年材料定位；实录记载反映朝廷选择，崇祯期须另以奏疏、地方及敌对政权材料交叉。 因而这里标示的是材料能够支持的行动与制度位置，而不是对动机、地方执行或普通人经验的无保留复原。

\eventanchor{ML-1500-002}{明·1500（弘治十三年）}
\paragraph{1500（弘治十三年）：清丈、库藏、差役与征解的本年记录显示财政监管压力，整饬命令不等于隐漏已被清除。} 这一节点应放回本章已经展开的权力、财政、地方社会与边地关系中理解。账本把它出现的条件概括为：财政征解、交通工程或货币支付的压力使相关措施进入政务。 随后可追踪的变化是：其法定安排不等于地方执行，后续须以地域材料追踪负担与成效。 可由《孝宗实录》弘治十三年条；《明会典》相关条目；地方志、黄册/契约/河工材料 互相核对；不过，桥接条目只标示可回查的年度问题与材料链；实录有中央选择性，崇祯期尤其需要奏疏、地方、朝鲜及满文材料交叉。 因而这里标示的是材料能够支持的行动与制度位置，而不是对动机、地方执行或普通人经验的无保留复原。

\eventanchor{ML-1500-003}{明·1500（弘治十三年）}
\paragraph{1500（弘治十三年）：嘉靖初继统礼制的本年文书构成大礼议过程锚点，礼制理由和政治效应不可简化为单一动机。（材料核查重点：诏令或奏议的形成与发受链）} 这一节点应放回本章已经展开的权力、财政、地方社会与边地关系中理解。账本把它出现的条件概括为：继承、官僚协调或皇权运作的压力使问题进入朝廷程序。 随后可追踪的变化是：它影响后续人事与政策议程；动机和派系标签不能由单一叙事裁定。 可由《孝宗实录》弘治十三年条；《明史·本纪》及相关列传；诏令、奏疏或会典版本 互相核对；不过，桥接条目只标示可回查的年度问题与材料链；实录有中央选择性，崇祯期尤其需要奏疏、地方、朝鲜及满文材料交叉。；本条特别核对诏令或奏议的形成与发受链。 因而这里标示的是材料能够支持的行动与制度位置，而不是对动机、地方执行或普通人经验的无保留复原。

\eventanchor{ML-1500-004}{明·1500（弘治十三年）}
\paragraph{1500（弘治十三年）：嘉靖初继统礼制的本年文书构成大礼议过程锚点，礼制理由和政治效应不可简化为单一动机。（材料核查重点：报称信息的来源、抄传与行政分类）} 这一节点应放回本章已经展开的权力、财政、地方社会与边地关系中理解。账本把它出现的条件概括为：继承、官僚协调或皇权运作的压力使问题进入朝廷程序。 随后可追踪的变化是：它影响后续人事与政策议程；动机和派系标签不能由单一叙事裁定。 可由《孝宗实录》弘治十三年条；《明史·本纪》及相关列传；诏令、奏疏或会典版本 互相核对；不过，桥接条目只标示可回查的年度问题与材料链；实录有中央选择性，崇祯期尤其需要奏疏、地方、朝鲜及满文材料交叉。；本条特别核对报称信息的来源、抄传与行政分类。 因而这里标示的是材料能够支持的行动与制度位置，而不是对动机、地方执行或普通人经验的无保留复原。

\eventanchor{ML-1500-005}{明·1500（弘治十三年）}
\paragraph{1500（弘治十三年）：嘉靖初继统礼制的本年文书构成大礼议过程锚点，礼制理由和政治效应不可简化为单一动机。（材料核查重点：同年地方材料所见的执行范围）} 这一节点应放回本章已经展开的权力、财政、地方社会与边地关系中理解。账本把它出现的条件概括为：继承、官僚协调或皇权运作的压力使问题进入朝廷程序。 随后可追踪的变化是：它影响后续人事与政策议程；动机和派系标签不能由单一叙事裁定。 可由《孝宗实录》弘治十三年条；《明史·本纪》及相关列传；诏令、奏疏或会典版本 互相核对；不过，桥接条目只标示可回查的年度问题与材料链；实录有中央选择性，崇祯期尤其需要奏疏、地方、朝鲜及满文材料交叉。；本条特别核对同年地方材料所见的执行范围。 因而这里标示的是材料能够支持的行动与制度位置，而不是对动机、地方执行或普通人经验的无保留复原。

\eventanchor{ML-1501-001}{明·1501（弘治十四年）}
\paragraph{1501（弘治十四年）：正德朝内廷、人事与巡幸相关文书构成本年中枢政治线索，后出道德化叙事须与原奏疏分开。（材料核查重点：诏令或奏议的形成与发受链）} 这一节点应放回本章已经展开的权力、财政、地方社会与边地关系中理解。账本把它出现的条件概括为：继承、官僚协调或皇权运作的压力使问题进入朝廷程序。 随后可追踪的变化是：它影响后续人事与政策议程；动机和派系标签不能由单一叙事裁定。 可由《孝宗实录》弘治十四年条；《明史·本纪》及相关列传；诏令、奏疏或会典版本 互相核对；不过，桥接条目只标示可回查的年度问题与材料链；实录有中央选择性，崇祯期尤其需要奏疏、地方、朝鲜及满文材料交叉。；本条特别核对诏令或奏议的形成与发受链。 因而这里标示的是材料能够支持的行动与制度位置，而不是对动机、地方执行或普通人经验的无保留复原。

\eventanchor{ML-1501-002}{明·1501（弘治十四年）}
\paragraph{1501（弘治十四年）：东南港口、日本使团或葡萄牙接触的本年材料为海上交往留档，朝贡、私贸与冲突不可混称。} 这一节点应放回本章已经展开的权力、财政、地方社会与边地关系中理解。账本把它出现的条件概括为：朝贡名义、边境安全与港口贸易的利益使交涉持续发生。 随后可追踪的变化是：它为后续册封、互市或海防安排留下文书与跨政权比较线索。 可由《孝宗实录》弘治十四年条；《明史·外国传》；《朝鲜王朝实录》、琉球《历代宝案》或海外档案 互相核对；不过，桥接条目只标示可回查的年度问题与材料链；实录有中央选择性，崇祯期尤其需要奏疏、地方、朝鲜及满文材料交叉。 因而这里标示的是材料能够支持的行动与制度位置，而不是对动机、地方执行或普通人经验的无保留复原。

\eventanchor{ML-1501-003}{明·1501（弘治十四年）}
\paragraph{1501（弘治十四年）：正德朝内廷、人事与巡幸相关文书构成本年中枢政治线索，后出道德化叙事须与原奏疏分开。（材料核查重点：报称信息的来源、抄传与行政分类）} 这一节点应放回本章已经展开的权力、财政、地方社会与边地关系中理解。账本把它出现的条件概括为：继承、官僚协调或皇权运作的压力使问题进入朝廷程序。 随后可追踪的变化是：它影响后续人事与政策议程；动机和派系标签不能由单一叙事裁定。 可由《孝宗实录》弘治十四年条；《明史·本纪》及相关列传；诏令、奏疏或会典版本 互相核对；不过，桥接条目只标示可回查的年度问题与材料链；实录有中央选择性，崇祯期尤其需要奏疏、地方、朝鲜及满文材料交叉。；本条特别核对报称信息的来源、抄传与行政分类。 因而这里标示的是材料能够支持的行动与制度位置，而不是对动机、地方执行或普通人经验的无保留复原。

\eventanchor{ML-1501-004}{明·1501（弘治十四年）}
\paragraph{1501（弘治十四年）：正德朝内廷、人事与巡幸相关文书构成本年中枢政治线索，后出道德化叙事须与原奏疏分开。（材料核查重点：同年地方材料所见的执行范围）} 这一节点应放回本章已经展开的权力、财政、地方社会与边地关系中理解。账本把它出现的条件概括为：继承、官僚协调或皇权运作的压力使问题进入朝廷程序。 随后可追踪的变化是：它影响后续人事与政策议程；动机和派系标签不能由单一叙事裁定。 可由《孝宗实录》弘治十四年条；《明史·本纪》及相关列传；诏令、奏疏或会典版本 互相核对；不过，桥接条目只标示可回查的年度问题与材料链；实录有中央选择性，崇祯期尤其需要奏疏、地方、朝鲜及满文材料交叉。；本条特别核对同年地方材料所见的执行范围。 因而这里标示的是材料能够支持的行动与制度位置，而不是对动机、地方执行或普通人经验的无保留复原。

\eventanchor{ML-1501-005}{明·1501（弘治十四年）}
\paragraph{1501（弘治十四年）：正德朝内廷、人事与巡幸相关文书构成本年中枢政治线索，后出道德化叙事须与原奏疏分开。（材料核查重点：原始日期、机构与地域之间的对应关系）} 这一节点应放回本章已经展开的权力、财政、地方社会与边地关系中理解。账本把它出现的条件概括为：继承、官僚协调或皇权运作的压力使问题进入朝廷程序。 随后可追踪的变化是：它影响后续人事与政策议程；动机和派系标签不能由单一叙事裁定。 可由《孝宗实录》弘治十四年条；《明史·本纪》及相关列传；诏令、奏疏或会典版本 互相核对；不过，桥接条目只标示可回查的年度问题与材料链；实录有中央选择性，崇祯期尤其需要奏疏、地方、朝鲜及满文材料交叉。；本条特别核对原始日期、机构与地域之间的对应关系。 因而这里标示的是材料能够支持的行动与制度位置，而不是对动机、地方执行或普通人经验的无保留复原。

\eventanchor{ML-1502-001}{明·1502（弘治十五年）}
\paragraph{1502（弘治十五年）：贵州彝族叛乱被镇压} 这一节点应放回本章已经展开的权力、财政、地方社会与边地关系中理解。账本把它出现的条件概括为：继承、官僚协调或皇权运作的压力使问题进入朝廷程序。 随后可追踪的变化是：它影响后续人事与政策议程；动机和派系标签不能由单一叙事裁定。 可由《孝宗实录》弘治十五年条；《明史·本纪》及相关列传；诏令、奏疏或会典版本 互相核对；不过，译写候选以编年材料定位；实录记载反映朝廷选择，崇祯期须另以奏疏、地方及敌对政权材料交叉。 因而这里标示的是材料能够支持的行动与制度位置，而不是对动机、地方执行或普通人经验的无保留复原。

\eventanchor{ML-1502-002}{明·1502（弘治十五年）}
\paragraph{1502（弘治十五年）：地方反抗、边镇调兵和宗藩危机的本年军政材料标记跨区域动员，军数与战果需要复核。（材料核查重点：诏令或奏议的形成与发受链）} 这一节点应放回本章已经展开的权力、财政、地方社会与边地关系中理解。账本把它出现的条件概括为：前期边防、地方武装或跨区域资源调动的压力推动了该行动。 随后可追踪的变化是：它改变了后续调兵、补给或地方秩序的条件；战果和伤亡须分开核对。 可由《孝宗实录》弘治十五年条；《明史·兵志》及相关列传；边镇、地方志或对方材料 互相核对；不过，桥接条目只标示可回查的年度问题与材料链；实录有中央选择性，崇祯期尤其需要奏疏、地方、朝鲜及满文材料交叉。；本条特别核对诏令或奏议的形成与发受链。 因而这里标示的是材料能够支持的行动与制度位置，而不是对动机、地方执行或普通人经验的无保留复原。

\eventanchor{ML-1502-003}{明·1502（弘治十五年）}
\paragraph{1502（弘治十五年）：地方反抗、边镇调兵和宗藩危机的本年军政材料标记跨区域动员，军数与战果需要复核。（材料核查重点：报称信息的来源、抄传与行政分类）} 这一节点应放回本章已经展开的权力、财政、地方社会与边地关系中理解。账本把它出现的条件概括为：前期边防、地方武装或跨区域资源调动的压力推动了该行动。 随后可追踪的变化是：它改变了后续调兵、补给或地方秩序的条件；战果和伤亡须分开核对。 可由《孝宗实录》弘治十五年条；《明史·兵志》及相关列传；边镇、地方志或对方材料 互相核对；不过，桥接条目只标示可回查的年度问题与材料链；实录有中央选择性，崇祯期尤其需要奏疏、地方、朝鲜及满文材料交叉。；本条特别核对报称信息的来源、抄传与行政分类。 因而这里标示的是材料能够支持的行动与制度位置，而不是对动机、地方执行或普通人经验的无保留复原。

\eventanchor{ML-1502-004}{明·1502（弘治十五年）}
\paragraph{1502（弘治十五年）：地方反抗、边镇调兵和宗藩危机的本年军政材料标记跨区域动员，军数与战果需要复核。（材料核查重点：同年地方材料所见的执行范围）} 这一节点应放回本章已经展开的权力、财政、地方社会与边地关系中理解。账本把它出现的条件概括为：前期边防、地方武装或跨区域资源调动的压力推动了该行动。 随后可追踪的变化是：它改变了后续调兵、补给或地方秩序的条件；战果和伤亡须分开核对。 可由《孝宗实录》弘治十五年条；《明史·兵志》及相关列传；边镇、地方志或对方材料 互相核对；不过，桥接条目只标示可回查的年度问题与材料链；实录有中央选择性，崇祯期尤其需要奏疏、地方、朝鲜及满文材料交叉。；本条特别核对同年地方材料所见的执行范围。 因而这里标示的是材料能够支持的行动与制度位置，而不是对动机、地方执行或普通人经验的无保留复原。

\eventanchor{MM-1502-HUIDIAN-PRESENTED}{明·1502（弘治十五年）}
\paragraph{1502（弘治十五年）：弘治续修《大明会典》进呈，汇集当时中央典章} 这一节点应放回本章已经展开的权力、财政、地方社会与边地关系中理解。账本把它出现的条件概括为：明初会典已难覆盖制度变动，朝廷需要重新整理官制、礼仪和财赋规则。 随后可追踪的变化是：新会典成为官员援引旧制的重要文本，也将不同时期制度压缩为相对稳定的规范图景。 可由《孝宗实录》弘治十五年条；弘治《大明会典》序及目录；《明史·职官志》 互相核对；不过，进呈事与成书阶段可靠；会典汇编的是规范和沿革，不可据此推定各地机构都按同一方式运行。 因而这里标示的是材料能够支持的行动与制度位置，而不是对动机、地方执行或普通人经验的无保留复原。

\eventanchor{ML-1503-001}{明·1503（弘治十六年）}
\paragraph{1503（弘治十六年）：李叛军被镇压} 这一节点应放回本章已经展开的权力、财政、地方社会与边地关系中理解。账本把它出现的条件概括为：继承、官僚协调或皇权运作的压力使问题进入朝廷程序。 随后可追踪的变化是：它影响后续人事与政策议程；动机和派系标签不能由单一叙事裁定。 可由《孝宗实录》弘治十六年条；《明史·本纪》及相关列传；诏令、奏疏或会典版本 互相核对；不过，译写候选以编年材料定位；实录记载反映朝廷选择，崇祯期须另以奏疏、地方及敌对政权材料交叉。 因而这里标示的是材料能够支持的行动与制度位置，而不是对动机、地方执行或普通人经验的无保留复原。

\eventanchor{ML-1503-002}{明·1503（弘治十六年）}
\paragraph{1503（弘治十六年）：科举、讲学和官员文集的本年线索可用于追踪知识网络，主张与行政结果须分开。} 这一节点应放回本章已经展开的权力、财政、地方社会与边地关系中理解。账本把它出现的条件概括为：制度、教育或知识传播的需要推动了相关行动或文本形成。 随后可追踪的变化是：它提供制度和传播史的锚点，不能据此推断社会整体接受。 可由《孝宗实录》弘治十六年条；《明史·选举志》或《艺文志》；文集、碑刻或书目材料 互相核对；不过，桥接条目只标示可回查的年度问题与材料链；实录有中央选择性，崇祯期尤其需要奏疏、地方、朝鲜及满文材料交叉。 因而这里标示的是材料能够支持的行动与制度位置，而不是对动机、地方执行或普通人经验的无保留复原。

\eventanchor{ML-1503-003}{明·1503（弘治十六年）}
\paragraph{1503（弘治十六年）：清丈、库藏、差役与征解的本年记录显示财政监管压力，整饬命令不等于隐漏已被清除。（材料核查重点：诏令或奏议的形成与发受链）} 这一节点应放回本章已经展开的权力、财政、地方社会与边地关系中理解。账本把它出现的条件概括为：财政征解、交通工程或货币支付的压力使相关措施进入政务。 随后可追踪的变化是：其法定安排不等于地方执行，后续须以地域材料追踪负担与成效。 可由《孝宗实录》弘治十六年条；《明会典》相关条目；地方志、黄册/契约/河工材料 互相核对；不过，桥接条目只标示可回查的年度问题与材料链；实录有中央选择性，崇祯期尤其需要奏疏、地方、朝鲜及满文材料交叉。；本条特别核对诏令或奏议的形成与发受链。 因而这里标示的是材料能够支持的行动与制度位置，而不是对动机、地方执行或普通人经验的无保留复原。

\eventanchor{ML-1503-004}{明·1503（弘治十六年）}
\paragraph{1503（弘治十六年）：清丈、库藏、差役与征解的本年记录显示财政监管压力，整饬命令不等于隐漏已被清除。（材料核查重点：报称信息的来源、抄传与行政分类）} 这一节点应放回本章已经展开的权力、财政、地方社会与边地关系中理解。账本把它出现的条件概括为：财政征解、交通工程或货币支付的压力使相关措施进入政务。 随后可追踪的变化是：其法定安排不等于地方执行，后续须以地域材料追踪负担与成效。 可由《孝宗实录》弘治十六年条；《明会典》相关条目；地方志、黄册/契约/河工材料 互相核对；不过，桥接条目只标示可回查的年度问题与材料链；实录有中央选择性，崇祯期尤其需要奏疏、地方、朝鲜及满文材料交叉。；本条特别核对报称信息的来源、抄传与行政分类。 因而这里标示的是材料能够支持的行动与制度位置，而不是对动机、地方执行或普通人经验的无保留复原。

\subsection{1503年前后的材料线索}

\eventanchor{ML-1503-005}{明·1503（弘治十六年）}
\paragraph{1503（弘治十六年）：清丈、库藏、差役与征解的本年记录显示财政监管压力，整饬命令不等于隐漏已被清除。（材料核查重点：同年地方材料所见的执行范围）} 这一节点应放回本章已经展开的权力、财政、地方社会与边地关系中理解。账本把它出现的条件概括为：财政征解、交通工程或货币支付的压力使相关措施进入政务。 随后可追踪的变化是：其法定安排不等于地方执行，后续须以地域材料追踪负担与成效。 可由《孝宗实录》弘治十六年条；《明会典》相关条目；地方志、黄册/契约/河工材料 互相核对；不过，桥接条目只标示可回查的年度问题与材料链；实录有中央选择性，崇祯期尤其需要奏疏、地方、朝鲜及满文材料交叉。；本条特别核对同年地方材料所见的执行范围。 因而这里标示的是材料能够支持的行动与制度位置，而不是对动机、地方执行或普通人经验的无保留复原。

\eventanchor{ML-1504-001}{明·1504（弘治十七年）}
\paragraph{1504（弘治十七年）：大同被蒙古人袭击} 这一节点应放回本章已经展开的权力、财政、地方社会与边地关系中理解。账本把它出现的条件概括为：前期边防、地方武装或跨区域资源调动的压力推动了该行动。 随后可追踪的变化是：它改变了后续调兵、补给或地方秩序的条件；战果和伤亡须分开核对。 可由《孝宗实录》弘治十七年条；《明史·兵志》及相关列传；边镇、地方志或对方材料 互相核对；不过，译写候选以编年材料定位；实录记载反映朝廷选择，崇祯期须另以奏疏、地方及敌对政权材料交叉。 因而这里标示的是材料能够支持的行动与制度位置，而不是对动机、地方执行或普通人经验的无保留复原。

\eventanchor{ML-1504-002}{明·1504（弘治十七年）}
\paragraph{1504（弘治十七年）：嘉靖初继统礼制的本年文书构成大礼议过程锚点，礼制理由和政治效应不可简化为单一动机。} 这一节点应放回本章已经展开的权力、财政、地方社会与边地关系中理解。账本把它出现的条件概括为：继承、官僚协调或皇权运作的压力使问题进入朝廷程序。 随后可追踪的变化是：它影响后续人事与政策议程；动机和派系标签不能由单一叙事裁定。 可由《孝宗实录》弘治十七年条；《明史·本纪》及相关列传；诏令、奏疏或会典版本 互相核对；不过，桥接条目只标示可回查的年度问题与材料链；实录有中央选择性，崇祯期尤其需要奏疏、地方、朝鲜及满文材料交叉。 因而这里标示的是材料能够支持的行动与制度位置，而不是对动机、地方执行或普通人经验的无保留复原。

\eventanchor{ML-1504-003}{明·1504（弘治十七年）}
\paragraph{1504（弘治十七年）：东南港口、日本使团或葡萄牙接触的本年材料为海上交往留档，朝贡、私贸与冲突不可混称。（材料核查重点：诏令或奏议的形成与发受链）} 这一节点应放回本章已经展开的权力、财政、地方社会与边地关系中理解。账本把它出现的条件概括为：朝贡名义、边境安全与港口贸易的利益使交涉持续发生。 随后可追踪的变化是：它为后续册封、互市或海防安排留下文书与跨政权比较线索。 可由《孝宗实录》弘治十七年条；《明史·外国传》；《朝鲜王朝实录》、琉球《历代宝案》或海外档案 互相核对；不过，桥接条目只标示可回查的年度问题与材料链；实录有中央选择性，崇祯期尤其需要奏疏、地方、朝鲜及满文材料交叉。；本条特别核对诏令或奏议的形成与发受链。 因而这里标示的是材料能够支持的行动与制度位置，而不是对动机、地方执行或普通人经验的无保留复原。

\eventanchor{ML-1504-004}{明·1504（弘治十七年）}
\paragraph{1504（弘治十七年）：东南港口、日本使团或葡萄牙接触的本年材料为海上交往留档，朝贡、私贸与冲突不可混称。（材料核查重点：报称信息的来源、抄传与行政分类）} 这一节点应放回本章已经展开的权力、财政、地方社会与边地关系中理解。账本把它出现的条件概括为：朝贡名义、边境安全与港口贸易的利益使交涉持续发生。 随后可追踪的变化是：它为后续册封、互市或海防安排留下文书与跨政权比较线索。 可由《孝宗实录》弘治十七年条；《明史·外国传》；《朝鲜王朝实录》、琉球《历代宝案》或海外档案 互相核对；不过，桥接条目只标示可回查的年度问题与材料链；实录有中央选择性，崇祯期尤其需要奏疏、地方、朝鲜及满文材料交叉。；本条特别核对报称信息的来源、抄传与行政分类。 因而这里标示的是材料能够支持的行动与制度位置，而不是对动机、地方执行或普通人经验的无保留复原。

\eventanchor{ML-1504-005}{明·1504（弘治十七年）}
\paragraph{1504（弘治十七年）：东南港口、日本使团或葡萄牙接触的本年材料为海上交往留档，朝贡、私贸与冲突不可混称。（材料核查重点：同年地方材料所见的执行范围）} 这一节点应放回本章已经展开的权力、财政、地方社会与边地关系中理解。账本把它出现的条件概括为：朝贡名义、边境安全与港口贸易的利益使交涉持续发生。 随后可追踪的变化是：它为后续册封、互市或海防安排留下文书与跨政权比较线索。 可由《孝宗实录》弘治十七年条；《明史·外国传》；《朝鲜王朝实录》、琉球《历代宝案》或海外档案 互相核对；不过，桥接条目只标示可回查的年度问题与材料链；实录有中央选择性，崇祯期尤其需要奏疏、地方、朝鲜及满文材料交叉。；本条特别核对同年地方材料所见的执行范围。 因而这里标示的是材料能够支持的行动与制度位置，而不是对动机、地方执行或普通人经验的无保留复原。

\eventanchor{ML-1505-001}{明·1505（弘治十八年）六月8日}
\paragraph{1505（弘治十八年）六月8日：弘治皇帝去世} 这一节点应放回本章已经展开的权力、财政、地方社会与边地关系中理解。账本把它出现的条件概括为：继承、官僚协调或皇权运作的压力使问题进入朝廷程序。 随后可追踪的变化是：它影响后续人事与政策议程；动机和派系标签不能由单一叙事裁定。 可由《孝宗实录》弘治十八年条；《明史·本纪》及相关列传；诏令、奏疏或会典版本 互相核对；不过，译写候选以编年材料定位；实录记载反映朝廷选择，崇祯期须另以奏疏、地方及敌对政权材料交叉。 因而这里标示的是材料能够支持的行动与制度位置，而不是对动机、地方执行或普通人经验的无保留复原。

\eventanchor{ML-1505-002}{明·1505（弘治十八年）六月19日}
\paragraph{1505（弘治十八年）六月19日：朱厚照即位正德皇帝} 这一节点应放回本章已经展开的权力、财政、地方社会与边地关系中理解。账本把它出现的条件概括为：继承、官僚协调或皇权运作的压力使问题进入朝廷程序。 随后可追踪的变化是：它影响后续人事与政策议程；动机和派系标签不能由单一叙事裁定。 可由《孝宗实录》弘治十八年条；《明史·本纪》及相关列传；诏令、奏疏或会典版本 互相核对；不过，译写候选以编年材料定位；实录记载反映朝廷选择，崇祯期须另以奏疏、地方及敌对政权材料交叉。 因而这里标示的是材料能够支持的行动与制度位置，而不是对动机、地方执行或普通人经验的无保留复原。

\eventanchor{ML-1505-003}{明·1505（弘治十八年）}
\paragraph{1505（弘治十八年）：正德皇帝开始使用太监担任军事和财政总管} 这一节点应放回本章已经展开的权力、财政、地方社会与边地关系中理解。账本把它出现的条件概括为：前期边防、地方武装或跨区域资源调动的压力推动了该行动。 随后可追踪的变化是：它改变了后续调兵、补给或地方秩序的条件；战果和伤亡须分开核对。 可由《孝宗实录》弘治十八年条；《明史·兵志》及相关列传；边镇、地方志或对方材料 互相核对；不过，译写候选以编年材料定位；实录记载反映朝廷选择，崇祯期须另以奏疏、地方及敌对政权材料交叉。 因而这里标示的是材料能够支持的行动与制度位置，而不是对动机、地方执行或普通人经验的无保留复原。

\eventanchor{ML-1505-004}{明·1505（弘治十八年）}
\paragraph{1505（弘治十八年）：科举、讲学和官员文集的本年线索可用于追踪知识网络，主张与行政结果须分开。} 这一节点应放回本章已经展开的权力、财政、地方社会与边地关系中理解。账本把它出现的条件概括为：制度、教育或知识传播的需要推动了相关行动或文本形成。 随后可追踪的变化是：它提供制度和传播史的锚点，不能据此推断社会整体接受。 可由《孝宗实录》弘治十八年条；《明史·选举志》或《艺文志》；文集、碑刻或书目材料 互相核对；不过，桥接条目只标示可回查的年度问题与材料链；实录有中央选择性，崇祯期尤其需要奏疏、地方、朝鲜及满文材料交叉。 因而这里标示的是材料能够支持的行动与制度位置，而不是对动机、地方执行或普通人经验的无保留复原。

\eventanchor{MM-1505-WUZONG-ACCESSION}{明·1505（弘治十八年）五月}
\paragraph{1505（弘治十八年）五月：孝宗去世，皇太子朱厚照即位，是为武宗} 这一节点应放回本章已经展开的权力、财政、地方社会与边地关系中理解。账本把它出现的条件概括为：朱厚照已为太子，皇位承接不存在景泰末年那样的名分竞争。 随后可追踪的变化是：少年皇帝即位改变了内廷近侍与文官大臣接近皇权的条件，刘瑾等人迅速上升。 可由《孝宗实录》弘治十八年五月条；《武宗实录》正德元年即位条；《明史·武宗本纪》 互相核对；不过，去世和即位可靠；临终政务安排仅能反映官修记录所保留的礼制程序。 因而这里标示的是材料能够支持的行动与制度位置，而不是对动机、地方执行或普通人经验的无保留复原。

\eventanchor{ML-1506-001}{明·1506（正德元年）五月}
\paragraph{1506（正德元年）五月：税务部责令调查收入不足问题} 这一节点应放回本章已经展开的权力、财政、地方社会与边地关系中理解。账本把它出现的条件概括为：继承、官僚协调或皇权运作的压力使问题进入朝廷程序。 随后可追踪的变化是：它影响后续人事与政策议程；动机和派系标签不能由单一叙事裁定。 可由《武宗实录》正德元年条；《明史·本纪》及相关列传；诏令、奏疏或会典版本 互相核对；不过，译写候选以编年材料定位；实录记载反映朝廷选择，崇祯期须另以奏疏、地方及敌对政权材料交叉。 因而这里标示的是材料能够支持的行动与制度位置，而不是对动机、地方执行或普通人经验的无保留复原。

\eventanchor{ML-1506-002}{明·1506（正德元年）七月}
\paragraph{1506（正德元年）七月：户部尚书韩文抱怨皇帝用大臣府库支出} 这一节点应放回本章已经展开的权力、财政、地方社会与边地关系中理解。账本把它出现的条件概括为：继承、官僚协调或皇权运作的压力使问题进入朝廷程序。 随后可追踪的变化是：它影响后续人事与政策议程；动机和派系标签不能由单一叙事裁定。 可由《武宗实录》正德元年条；《明史·本纪》及相关列传；诏令、奏疏或会典版本 互相核对；不过，译写候选以编年材料定位；实录记载反映朝廷选择，崇祯期须另以奏疏、地方及敌对政权材料交叉。 因而这里标示的是材料能够支持的行动与制度位置，而不是对动机、地方执行或普通人经验的无保留复原。

\eventanchor{ML-1506-003}{明·1506（正德元年）十月28日}
\paragraph{1506（正德元年）十月28日：户部尚书奏请皇帝将他亲手的太监全部处死，皇帝不肯，大臣们都辞职了。} 这一节点应放回本章已经展开的权力、财政、地方社会与边地关系中理解。账本把它出现的条件概括为：继承、官僚协调或皇权运作的压力使问题进入朝廷程序。 随后可追踪的变化是：它影响后续人事与政策议程；动机和派系标签不能由单一叙事裁定。 可由《武宗实录》正德元年条；《明史·本纪》及相关列传；诏令、奏疏或会典版本 互相核对；不过，译写候选以编年材料定位；实录记载反映朝廷选择，崇祯期须另以奏疏、地方及敌对政权材料交叉。 因而这里标示的是材料能够支持的行动与制度位置，而不是对动机、地方执行或普通人经验的无保留复原。
